\begin{document}
\hypersetup{colorlinks, linkcolor=blue}
\maketitle

\section{\textbf{Appointments}}
\cventry{2019- present}{Postdoctoral scholar, Sali Lab, Department of Bioengineering and Therapeutic Sciences}{University of California San Francisco}{San Francisco}{CA}{}

\section{\textbf{Education}}
\cventry{2013- 2018}{Ph.D, Chemical Engineering}{University of California Santa Barbara}{Santa Barbara}{CA}{}

\cventry{2014- 2016}{Graduate emphasis in Computational Science and Engineering}{University of California Santa Barbara}{Santa Barbara}{CA}{(Specialization in numerical methods for differential equations, and parallel computing)}

\cventry{2008- 2013}{Masters and Bachelors of Technology (Hons.) integrated dual degree, Chemical Engineering}{Indian Institute of Technology Kharagpur}{Kharagpur}{WB, India}{}

\section{\textbf{Research Interests}}
\cvline{}{Integrative modeling of protein complexes, chemical crosslinking, protein folding, probabilistic machine learning and Bayesian inference, whole cell modeling.}

\section{\textbf{Skills}}
\cvline{Softwares}{MATLAB, Mathematica, LAMMPS, GROMACS, AMBER, IMP, VMD, PyMol, UCSF Chimera}
\cvline{Frameworks}{Tensorflow-2,  PyMC3}
\cvline{Languages}{Python (\textasciitilde 10000 lines), Fortran-90 (\textasciitilde 1000 lines), C++ (\textasciitilde 1800 lines), Bash, TCL}
\cvline{Parallel-platforms}{MPI, OpenMP}

\section{\textbf{Research Experience}}
\cvline{\textbf{UCSF}}{Postdoc advisor: Andrej Sali}
\cvlistitem{\textbf{Integrative structural biology:} \newline Integrative structure determination of biologically relevant macromolecular complexes using experimental chemical crosslinking data. Work involves both developing Bayesian inference based methods for improved structure determination and applying such methods to study chromosomal protein complexes and DNA helicases.}

\vspace{0.1in}
\cvlistitem{\textbf{Whole cell models:} \newline As a part of the Pancreatic Beta-Cell Consortium (collaborating across USC, UCLA, UCSF and SanghaiTech), I am working towards coupling traditional models of different parts, aspects, length and time-scales of beta cell biology, to build a \textit{meta-model} of the whole-cell. Work involves developing computational toolchains for translating several different models and data like macroscopic pharmacokinetics, mesoscopic spatio-temporal Brownian dynamics simulations, microscopic molecular dynamics simulations of protein-protein interactions, biochemical pathways, soft-X-ray tomograms of insulin vesicles and other proteomic data into probabilistic graphical models as a first step to couple such seemingly disparate models.}
\vspace{0.1in}

\cvline{\textbf{UCSB}}{PhD advisor: M. Scott Shell}
\cvlistitem{\textbf{Thesis: Next-Generation Coarse-Grained Models for Molecular Dynamics Simulations of Fluid Phase Equilibria and Protein Biophysics Using the Relative Entropy} \newline Introduced the local density potential for structurally and thermodynamically accurate coarse-grained molecular dynamics (MD) simulations of liquid-liquid phase separation, which are relatively rare in the chemical engineering literature. \newline \newline Also developed protein backbone models for (a) template-free folding of 200+ residue coarse-grained protein domains using MD simulations and, (b) studying self-assembly in amyloidogenic peptides (as seen in neurodegenerative diseases like Alzheimer's, ALS, etc).}
\vspace{0.1in}

\cvline{\textbf{IIT Kharagpur}}{Undergraduate thesis advisor: Saikat Chakraborty}
\cvlistitem{Reactive-diffusive modeling of human respiration and its application in characterizing the patho-physiology of methemoglobin induced anemia. (Awarded best thesis)}

\section{\textbf{Publications}}
\cvline{Google Scholar}{\url{https://scholar.google.com/citations?hl=en&user=MastYM0AAAAJ}}
\cvlistitem{T. Sanyal, J. Mittal and M.S. Shell,\textit{ A hybrid, bottom up, structurally-accurate, G$\bar{o}$-like coarse-grained protein model}. Journal of Chemical Physics (2019) 151, 044111 (\textit{Editor's pick article, 2019})}
\vspace{0.1in}
\cvlistitem{D. Rosenberger, T.Sanyal, M.S. Shell and N.F.A van der Vegt, \textit{Transferability of local density-assisted implicit solvation models for homogeneous fluid mixtures.} Journal of Chemical Theory and Computation (2019) 15, 2881-2895}
\vspace{0.1in}
\cvlistitem{M.P. Howard, W.F. Reinhart, T. Sanyal, M.S. Shell, A. Nikoubashman and A.Z. Panagiotopoulos, \textit{Evaporation-induced assembly of colloidal crystals.} Journal of Chemical Physics (2018) 149(\textbf{9})}
\vspace{0.1in}
\cvlistitem{T. Sanyal and M.S. Shell, \textit{Transferable coarse-grained models of liquid-liquid equilibrium using local density potentials optimized with the relative entropy.} Journal of Physical Chemistry-B (2018) 122 (\textbf{21}), 5678-5693}
\vspace{0.1in}
\cvlistitem{T. Sanyal and M.S. Shell, \textit{Coarse-grained models using local-density potentials optimized with the relative entropy: Application to implicit solvation.} Journal of Chemical Physics (2016) \textbf{145}, 034109}
\vspace{0.1in}
\cvlistitem{T. Sanyal and S. Chakraborty, \textit{Multiscale analysis of hypoxemia in methemoglobin-anemia.} Mathematical Biosciences (2013) \textbf{241}, 167-180}
\vspace{0.1in}
\cvlistitem{T. Sanyal and S. Chakraborty, \textit{Multiscale analysis of simultaneous uptake of two interacting gases in the human lungs and its application to methemoglobin anemia.} Computers and Chemical Engineering (2013) \textbf{59}, 226-242}
\vspace{0.1in}

\section{\textbf{Outreach Activities}}
\cvline{2014-2016}{\textbf{Co-founder and co-organizer, ChE Graduate Simulation Seminar Series, UCSB}}
\cvlistitem{Conceptualized and started a seminar series with colleagues to address the unfulfilled need for a common collaborative and social platform for theory and computation researchers on the UCSB campus.}
\cvlistitem{Developed extensive public outreach and attracted funding from ChE, Materials and Computer Science faculty, with participation from over 10 STEM fields and special attendances and keynotes from industry and national labs}
\end{document}
